\chapter{Παραδείγματα Χρήσης}
\label{chap:usage}

Το παρόν κεφάλαιο αποτελεί πρακτικό οδηγό χρήσης του συστήματος: παρουσιάζει βήμα προς βήμα την εγκατάσταση, την επίλυση μιας διαταραγμένης κατάστασης, την επιλογή επιλυτή, την αναγνώριση απόστασης από τη λύση και την εκτέλεση των μετρήσεων. Στόχος του είναι να μπορεί ένας αναγνώστης να χρησιμοποιήσει το αποτέλεσμα της εργασίας χωρίς να ανατρέξει στα τεχνικά κεφάλαια. Η εσωτερική λειτουργία των επιλυτών περιγράφεται στο Κεφάλαιο~\ref{chap:implementation}, τα πειραματικά αποτελέσματα στο Κεφάλαιο~\ref{chap:experiments}, ενώ η εξαντλητική αναφορά εντολών και σημαιών δίνεται στο Παράρτημα~\ref{app:cli}. Όλες οι εντολές που ακολουθούν αντιστοιχούν στην πραγματική διεπαφή γραμμής εντολών της υλοποίησης.

\section{Προαπαιτούμενα και εγκατάσταση}
\label{sec:usage-install}

Η υλοποίηση απαιτεί Python έκδοσης 3.10 ή νεότερης. Η εγκατάσταση γίνεται σε εικονικό περιβάλλον (\eng{virtual environment}) με εγκατάσταση του πακέτου σε επεξεργάσιμη μορφή (\eng{editable mode}):

\begin{verbatim}
python -m pip install -e ".[dev]"
rubik-optimal --help
\end{verbatim}

Η εγκατάσταση καταχωρίζει την εντολή κονσόλας \code{rubik-optimal}· ισοδύναμα, η διεπαφή καλείται ως \code{python -m rubik_optimal.cli}. Η προεπιλεγμένη εγκατάσταση δεν εισάγει κανέναν εξωτερικό επιλυτή: οι υλοποιήσεις Thistlethwaite, Kociemba δύο φάσεων και Korf/IDA* λειτουργούν χωρίς προαιρετικές εξαρτήσεις. Το προαιρετικό πρόσθετο \code{baselines} εγκαθιστά το εξωτερικό πακέτο \code{kociemba}, το οποίο χρησιμοποιείται μόνο ως βάση σύγκρισης.

Οι εγγενείς (\eng{native}) διαδρομές ακριβούς αναζήτησης απαιτούν διαθέσιμο μεταγλωττιστή C/C++· τα αντίστοιχα σενάρια παραγωγής μεταγλωττίζουν αυτόματα τους εγγενείς γεννήτορες κατά την πρώτη εκτέλεση. Πριν από τη χρήση των επιλυτών που στηρίζονται σε πίνακες, πρέπει να παραχθούν οι πίνακες κινήσεων (\eng{move tables}) και οι πίνακες αποκοπής (\eng{pruning tables}) των συντεταγμένων:

\begin{verbatim}
rubik-optimal tables --profile thesis --seed 2026
\end{verbatim}

Οι μεγαλύτερες δομές, δηλαδή οι πίνακες προτύπων (\eng{pattern databases}, PDB) γωνιών και ακμών και οι πίνακες H48, παράγονται από τα σενάρια \path{scripts/generate_corner_pdb.py}, \path{scripts/generate_edge_pdb.py} και \path{scripts/generate_h48_tables.py} αντίστοιχα. Η πλήρης αλληλουχία παραγωγής όλων των τεχνουργημάτων (\eng{artifacts}) περιγράφεται στο Παράρτημα~\ref{app:repro}.

\section{Επίλυση μιας διαταραγμένης κατάστασης}
\label{sec:usage-solve}

Η διεπαφή δέχεται την αρχική κατάσταση του κύβου 3×3×3 σε δύο μορφές. Η πρώτη είναι ακολουθία κινήσεων από τη λυμένη κατάσταση, στη σημειογραφία που ορίζεται στο Κεφάλαιο~\ref{chap:model}, π.χ. \code{"R U R' U'"}. Η δεύτερη είναι συμβολοσειρά εδρών (\eng{facelet string}) 54 χαρακτήρων πάνω στο αλφάβητο \code{URFDLB}, με τις έδρες σε σειρά U, R, F, D, L, B και κάθε χρώμα να εμφανίζεται ακριβώς εννέα φορές. Η κενή είσοδος ή η λέξη \code{solved} δηλώνουν τη λυμένη κατάσταση. Αναπαραγώγιμες διαταράξεις (\eng{scrambles}) παράγονται ντετερμινιστικά από σπόρο:

\begin{verbatim}
rubik-optimal scramble --length 20 --seed 2026
rubik-optimal facelets --length 20 --seed 2026
\end{verbatim}

Η πρώτη εντολή τυπώνει ακολουθία 20 κινήσεων και η δεύτερη την αντίστοιχη έγκυρη συμβολοσειρά εδρών· με την επιλογή \code{--json} η δεύτερη επιστρέφει επιπλέον τη διατάραξη και τα μεταδεδομένα της. Η επίλυση γίνεται με την εντολή \code{solve}:

\begin{verbatim}
rubik-optimal solve "R U R' U'" --timeout 5
\end{verbatim}

Πριν από κάθε αναζήτηση, η είσοδος ελέγχεται για φυσική εγκυρότητα (ισοτιμίες μεταθέσεων, προσανατολισμοί γωνιών και ακμών)· μη επιλύσιμες καταστάσεις απορρίπτονται με επεξηγηματικό μήνυμα πριν κληθεί οποιοσδήποτε επιλυτής. Το αποτέλεσμα τυπώνεται ως αντικείμενο JSON με τα εξής βασικά πεδία: \code{solver_name} (ο επιλυτής που παρήγαγε τη λύση), \code{solution_moves} και \code{solution_length} (η λύση και το μήκος της στη μετρική μισών στροφών, \eng{half-turn metric}, HTM), \code{status}, \code{is_verified}, \code{runtime_seconds} και \code{notes}. Το πεδίο \code{status} διακρίνει την ισχύ του αποτελέσματος: \code{exact} σημαίνει ότι η αναζήτηση ολοκληρώθηκε και η λύση είναι αποδεδειγμένα βέλτιστη, ενώ \code{non_exact} σημαίνει ότι η λύση επαληθεύτηκε αλλά δεν αποδεικνύεται ελάχιστη. Η ετικέτα \code{lower_bound} δηλώνει αποδεδειγμένο κάτω φράγμα χωρίς λύση και η \code{timeout} δηλώνει ότι η αναζήτηση σταμάτησε από όριο χρόνου, βάθους ή κόμβων. Το πεδίο \code{is_verified} είναι ανεξάρτητο από το \code{status}: δηλώνει ότι ο ανεξάρτητος επαληθευτής εφάρμοσε τη λύση στην αρχική κατάσταση και επιβεβαίωσε ότι καταλήγει στη λυμένη (Κεφάλαιο~\ref{chap:validation}).

Μια προτεινόμενη λύση επαληθεύεται και χειροκίνητα με την εντολή \code{verify}, η οποία δέχεται την αρχική κατάσταση και την υποψήφια ακολουθία. Για παράδειγμα, η αντίστροφη ακολουθία μιας διατάραξης τη λύνει πάντα:

\begin{verbatim}
rubik-optimal verify "R U R' U'" "U R U' R'"
\end{verbatim}

Η έξοδος περιέχει τα πεδία \code{ok}, \code{move_count} και \code{message}, και ο κωδικός εξόδου της εντολής είναι μηδενικός μόνο όταν η ακολουθία πράγματι λύνει την κατάσταση.

\section{Επιλογή επιλυτή και προφίλ}
\label{sec:usage-solvers}

Η επιλογή \code{--solver} της εντολής \code{solve} καθορίζει τη διαδρομή επίλυσης. Η προεπιλογή \code{auto} δοκιμάζει διαδοχικά την εγγενή διαδρομή δύο φάσεων Kociemba, τον εξωτερικό προσαρμογέα \code{kociemba} (εφόσον είναι εγκατεστημένος) και τη διαδρομή Thistlethwaite, και επιστρέφει την πρώτη επαληθευμένη λύση. Οι κυριότερες ρητές επιλογές είναι οι εξής:

\begin{itemize}
\item \code{native-kociemba}: η εγγενής υλοποίηση δύο φάσεων του Kociemba, με προεπιλεγμένα όρια βάθους ίσα με τις φασικές διαμέτρους (12 και 18). Λύνει τυπικές τυχαίες διαταράξεις 18--20 κινήσεων και επιστρέφει \code{non_exact}, καθώς η μέθοδος δεν αποδεικνύει βελτιστότητα.
\item \code{thistlethwaite}: η τετραφασική κάθοδος υποομάδων με ακριβείς φασικούς πίνακες αποστάσεων. Τερματίζει πάντα και επιστρέφει \code{non_exact} για κάθε επαληθευμένη λύση.
\item \code{korf}: η εκπαιδευτική IDA* σε Python, με παραδεκτό κάτω φράγμα από τους παραγόμενους πίνακες. Είναι η μόνη διαδρομή που τηρεί κατά γράμμα τα \code{--max-depth} (προεπιλογή 10) και \code{--node-limit} (προεπιλογή 500\,000)· κατάλληλη για ρηχές καταστάσεις.
\item \code{optimal-native}: ο εγγενής επιλυτής IDA* πλήρους κύβου με τον γωνιακό και τους οκτώ εξακμικούς πίνακες προτύπων. Όταν η αναζήτηση ολοκληρώνεται, το αποτέλεσμα είναι \code{exact} και η λύση βέλτιστη στην HTM.
\item \code{h48-native} και \code{h48-oracle}: η διαδρομή ακριβούς αναζήτησης H48 πάνω στο ενσωματωμένο υποσύστημα (\eng{backend}) \code{nissy-core} \cite{trontoNissyCoreH48}. Η πρώτη χρησιμοποιεί από προεπιλογή τον μικρό πίνακα \code{h48h0}, ενώ η δεύτερη επιλέγει τον μεγαλύτερο διαθέσιμο παραγόμενο πίνακα (στο προφίλ της εργασίας τον \code{h48h7}) και ενεργοποιεί αυτόματα τη χρήση των έμπιστων μεταδεδομένων του.
\item \code{kociemba} ή \code{adapter}: ο εξωτερικός προσαρμογέας του πακέτου \code{kociemba}, ως πρακτική βάση σύγκρισης με αποτελέσματα \code{non_exact}.
\end{itemize}

Κοινές παράμετροι για όλες τις διαδρομές είναι το \code{--timeout} (προεπιλογή 5 δευτερόλεπτα) και το \code{--threads}. Σημειώνεται ότι οι εγγενείς ακριβείς διαδρομές αυξάνουν σιωπηρά το \code{--max-depth} σε τουλάχιστον 20, ώστε καμία έγκυρη κατάσταση να μην αποκλείεται από τη διαμόρφωση. Για τον πίνακα \code{h48h7} διατίθενται επιπλέον οι επιλογές \code{--h48-trusted-table}, που παραλείπει την πλήρη σάρωση ελέγχου του πίνακα όταν αυτός έχει παραχθεί και ταυτοποιηθεί από μεταδεδομένα και αθροίσματα ελέγχου, και \code{--h48-preload-table}, που προθερμαίνει σειριακά τις σελίδες του πίνακα πριν από δύσκολες αναζητήσεις. Παράδειγμα κλήσης της ισχυρότερης διαδρομής σε διατάραξη 20 κινήσεων:

\begin{verbatim}
rubik-optimal solve \
  "$(rubik-optimal scramble --length 20 --seed 2026)" \
  --solver h48-oracle --timeout 90
\end{verbatim}

Η διεπαφή προσφέρει και συνδυαστικές διαδρομές χαρτοφυλακίου (\code{race-optimal}, \code{resident-race-optimal}, \code{universal-optimal}), οι οποίες εκτελούν παράλληλα ή κλιμακωτά περισσότερες ακριβείς μηχανές (\eng{backends})· οι πλήρεις επιλογές τους τεκμηριώνονται στο Παράρτημα~\ref{app:cli}.

\section{Αναγνώριση απόστασης από τη λύση}
\label{sec:usage-distance}

Ο δεύτερος στόχος της εργασίας, η αναγνώριση της απόστασης μιας κατάστασης από τη λύση, εξυπηρετείται από την εντολή \code{distance}. Η εντολή εκτελεί πρώτα ρηχή εξαντλητική αναζήτηση BFS μέχρι το βάθος \code{--bfs-depth} (προεπιλογή 5)· αν η κατάσταση βρεθεί, η απόσταση είναι αποδεδειγμένα ακριβής. Διαφορετικά, εκτελείται IDA* με το συνδυασμένο παραδεκτό κάτω φράγμα των παραγόμενων πινάκων μέχρι το βάθος \code{--ida-depth} (προεπιλογή 8), και αν ούτε αυτή ολοκληρωθεί, επιστρέφεται τεκμηριωμένο κάτω φράγμα:

\begin{verbatim}
rubik-optimal distance "R U R' U'"
\end{verbatim}

Το αποτέλεσμα είναι αντικείμενο JSON με πεδία \code{distance_value}, \code{kind}, \code{method}, \code{runtime_seconds}, \code{expanded_nodes} και \code{proof_notes}. Το πεδίο \code{kind} διακρίνει το είδος της απάντησης: \code{exact_distance} σημαίνει αποδεδειγμένη ακριβή απόσταση, \code{lower_bound} σημαίνει παραδεκτό κάτω φράγμα που αποκλείει μικρότερες αποστάσεις χωρίς να αποτελεί το ίδιο την απόσταση, \code{unknown_timeout} σημαίνει ότι η αναζήτηση δεν ολοκληρώθηκε εντός του ορίου και η τιμή πρέπει να διαβαστεί μόνο ως κάτω φράγμα, και \code{invalid_state} δηλώνει φυσικά μη έγκυρη είσοδο. Το πεδίο \code{method} καταγράφει ποια διαδρομή παρήγαγε την απάντηση, ώστε κάθε τιμή να είναι ιχνηλατήσιμη.

Για βαθύτερες καταστάσεις, η ακριβής αναγνώριση απαιτεί τις ισχυρότερες μηχανές. Η επιλογή \code{--optimal-native} χρησιμοποιεί τον εγγενή επιλυτή με τους πίνακες προτύπων, ενώ οι επιλογές \code{--h48-native} και \code{--h48-oracle} ενεργοποιούν τη διαδρομή H48 με είσοδο απευθείας την κατάσταση, χωρίς να απαιτείται η ακολουθία που την παρήγαγε:

\begin{verbatim}
rubik-optimal distance \
  "$(rubik-optimal facelets --length 20 --seed 2026)" \
  --h48-oracle --timeout 120
\end{verbatim}

Όταν η αναζήτηση H48 ολοκληρώνεται και η λύση επαληθεύεται, το \code{kind} είναι \code{exact_distance} και η τιμή είναι η ακριβής απόσταση της κατάστασης στην HTM. Η ίδια H48 διαδρομή απευθείας κατάστασης πιστοποίησε, μεταξύ άλλων, την ακριβή απόσταση 20 του superflip, όπως παρουσιάζεται στο Κεφάλαιο~\ref{chap:experiments}.

\section{Εκτέλεση δέσμης μετρήσεων και παραγωγή πινάκων}
\label{sec:usage-benchmarks}

Τα ποσοτικά αποτελέσματα της εργασίας παράγονται από το σενάριο δέσμης μετρήσεων με σταθερό προφίλ και σπόρο:

\begin{verbatim}
python scripts/run_benchmarks.py --profile thesis --seed 2026
\end{verbatim}

Το προφίλ \code{thesis} με σπόρο 2026 εκτελεί όλους τους επιλυτές πάνω στο σταθερό σώμα περιπτώσεων του Κεφαλαίου~\ref{chap:experiments}: τη λυμένη κατάσταση, ρηχές ντετερμινιστικές διαταράξεις βάθους 1 έως 8 και βαθύτερες ντετερμινιστικές τυχαίες περιπτώσεις βάθους 5, 10, 15 και 20. Τα παραγόμενα τεχνουργήματα αποθηκεύονται σε σταθερές διαδρομές: οι ωμές γραμμές αποτελεσμάτων στο \path{results/raw/benchmarks_seed_2026_thesis.jsonl} και η επεξεργασμένη σύνοψη με τον πίνακα CSV στο \path{results/processed/}. Η ισοδύναμη εντολή \code{rubik-optimal benchmark --profile thesis --seed 2026} εκτελεί την ίδια δέσμη· το προφίλ \code{quick} προορίζεται μόνο για γρήγορους δοκιμαστικούς ελέγχους και δεν παράγει τα αρχεία αναφοράς της εργασίας.

Μετά από κάθε εκτέλεση μετρήσεων, η συνέπεια των αποθηκευμένων αποτελεσμάτων ελέγχεται από τον επαληθευτή αποτελεσμάτων:

\begin{verbatim}
python scripts/verify_results.py
\end{verbatim}

Ο έλεγχος επιβεβαιώνει το σχήμα των γραμμών (Παράρτημα~\ref{app:schema}), επαναλαμβάνει την επαλήθευση κάθε δηλωμένης λύσης, ελέγχει τη συνέπεια των ρηχών ακριβών αποστάσεων με BFS, την ακεραιότητα των πινάκων προτύπων και την πληρότητα της κατανομής του κύβου 2×2×2. Τέλος, οι πίνακες LaTeX που εμφανίζονται στο Κεφάλαιο~\ref{chap:experiments} αναπαράγονται από τα αποθηκευμένα αποτελέσματα με το \code{python scripts/generate_figures.py} και γράφονται στον κατάλογο \path{thesis/tables/}. Πρακτικός κανόνας ονοματοδοσίας: όλα τα αρχεία αναφοράς της εργασίας περιέχουν στο όνομά τους το προφίλ \code{thesis} και τον σπόρο 2026, ώστε να μη συγχέονται με δοκιμαστικές εκτελέσεις.

Για επαναλαμβανόμενη ακριβή επίλυση πολλών καταστάσεων, η εντολή \code{oracle} δέχεται είσοδο γραμμή προς γραμμή από ορίσματα, αρχείο ή την πρότυπη είσοδο, και φορτώνει τον πίνακα \code{h48h7} μία φορά για όλο το σύνολο:

\begin{verbatim}
rubik-optimal oracle --input-file states.txt \
  --h48-trusted-table --timeout 300
\end{verbatim}

Η έξοδος είναι συγκεντρωτικό αντικείμενο JSON με μία γραμμή ανά είσοδο και συνολικά πεδία \code{all_exact}, \code{all_verified} και \code{batch_wall_seconds}· με την επιλογή \code{--jsonl} τυπώνεται μία γραμμή JSON ανά είσοδο, ενώ με την επιλογή \code{--stream} διατηρείται μόνιμα ενεργή εγγενής διεργασία και κάθε αποτέλεσμα τυπώνεται μόλις ολοκληρωθεί. Ο κωδικός εξόδου είναι μηδενικός μόνο όταν όλες οι γραμμές είναι \code{exact} και επαληθευμένες.

\section{Συνηθισμένα σφάλματα και όρια}
\label{sec:usage-errors}

Η συχνότερη παρανόηση αφορά τη σχέση των πεδίων \code{status} και \code{is_verified}. Η τιμή \code{is_verified=true} δεν σημαίνει βέλτιστη λύση: σημαίνει μόνο ότι η ακολουθία λύνει την κατάσταση. Η βελτιστότητα προκύπτει αποκλειστικά από την ετικέτα \code{exact}, η οποία αποδίδεται μόνο όταν ολοκληρωθεί παραδεκτή ή εξαντλητική αναζήτηση. Ο χρήστης πρέπει να διαβάζει και τα δύο πεδία.

Το προεπιλεγμένο όριο χρόνου των 5 δευτερολέπτων επαρκεί για τις πρακτικές διαδρομές και για ρηχές ακριβείς αναζητήσεις, όχι όμως για βαθιές καταστάσεις. Όταν το όριο εξαντληθεί, το αποτέλεσμα φέρει την ετικέτα \code{timeout}: δεν υπάρχει λύση στο αποτέλεσμα και το πεδίο \code{notes} εξηγεί τι συνέβη. Η ετικέτα αυτή καταγράφει μόνο τη μη ολοκλήρωση εντός του ορίου στο συγκεκριμένο μηχάνημα· δεν αποτελεί φράγμα απόστασης. Αντίστοιχα, μια απάντηση \code{lower_bound} αποκλείει μικρότερες αποστάσεις αλλά δεν δίνει λύση. Για δύσκολες μεμονωμένες καταστάσεις πρέπει να αυξάνεται ρητά το \code{--timeout}: ενδεικτικά, η πιστοποίηση του superflip μέσω της διαδρομής \code{h48h7} χρειάστηκε 51.813928 δευτερόλεπτα σε κλήση μίας διεργασίας (Κεφάλαιο~\ref{chap:experiments}).

Τα όρια μνήμης είναι εξίσου πρακτικά. Ο πίνακας \code{h48h7} έχει μέγεθος 3\,793\,842\,344 bytes και η πρώτη ψυχρή φόρτωσή του σε μηχάνημα 16~GiB είναι αισθητά ακριβή· για επαναλαμβανόμενη χρήση προτιμώνται η εντολή \code{oracle}, που φορτώνει τον πίνακα μία φορά, και οι επιλογές \code{--h48-trusted-table} και \code{--h48-preload-table}. Η σημαία \code{--h48-trusted-table} είναι ασφαλής μόνο για πίνακες που έχουν παραχθεί τοπικά και ταυτοποιηθεί από τα αποθηκευμένα μεταδεδομένα και αθροίσματα ελέγχου· σε κάθε άλλη περίπτωση πρέπει να επιτρέπεται ο πλήρης έλεγχος του πίνακα. Πίνακες μεγαλύτερων επιπέδων από τον \code{h48h7} δεν παράχθηκαν στο μηχάνημα αναφοράς και η εργασία δεν στηρίζει κανέναν ισχυρισμό σε αυτούς.

Τέλος, για τη διάγνωση σφαλμάτων εισόδου: συμβολοσειρά εδρών με λάθος μήκος, χαρακτήρες εκτός του \code{URFDLB}, λάθος πλήθος χρωμάτων ή φυσικά μη επιλύσιμη διάταξη απορρίπτεται με συγκεκριμένο μήνυμα πριν από οποιαδήποτε αναζήτηση. Οι εντολές \code{solve}, \code{verify} και \code{oracle} επιστρέφουν μη μηδενικό κωδικό εξόδου σε αποτυχία, ετικέτα \code{timeout} ή ανεπιτυχή επαλήθευση, ώστε η χρήση τους σε σενάρια κελύφους να είναι αξιόπιστη χωρίς ανάλυση του JSON.
